\documentclass[addpoints,10pt]{exam}

\usepackage{amsmath,amsthm,enumitem,wrapfig,amsfonts,mathtools}
\usepackage[mathscr]{euscript}
\usepackage[super]{nth}
\usepackage{dsfont}
\usepackage{xparse}
\usepackage{cancel}

\usepackage{geometry}
\usepackage[T1]{fontenc} % Use 8-bit encoding that has 256 glyphs
\renewcommand{\rmdefault}{ptm} %Change the Front Family from the default(cmr) to ptm(Times)
\usepackage{amsmath,amsfonts,amsthm,amssymb} % Math packages
\usepackage{bm}
%\usepackage{mathptmx}
\usepackage{graphicx}
\usepackage{sectsty} % Allows customizing section commands
% \allsectionsfont{\centering} % Make all sections centered, the default font and small caps
\usepackage{pgfplots}
\pgfplotsset{compat=1.18}
\usetikzlibrary{arrows.meta}
\usepackage{xcolor}
\definecolor{darkpastelgreen}{rgb}{0.01, 0.75, 0.24}
\definecolor{blue-violet}{rgb}{0.54, 0.17, 0.89}
%\usepackage{polylongdiv} %polynomial long division
% Custom problem environment
\usepackage{polynom}
\newcounter{cprob}
\newenvironment{cprob}[1]{%
    \setcounter{cprob}{#1}%
    \noindent\textbf{Problem \thecprob.}%
}{%
    \par\bigskip%
}

\theoremstyle{plain}
\newcommand{\theoremname}{Theorem}
\newtheorem{thm}{\protect\theoremname}
  \theoremstyle{definition}
  \newtheorem{prob}[thm]{Problem}
  \newtheorem*{problem*}{Open Problem}
  \theoremstyle{plain}
  \newtheorem{conjecture}[thm]{Conjecture}
  \theoremstyle{plain}
  \newtheorem{lem}[thm]{Lemma}
  \newtheorem*{lem*}{Lemma}
  \newtheorem{obs}[thm]{Observation}
  \newtheorem{cor}[thm]{Corollary}
  \theoremstyle{definition}
\newtheorem{definition}[thm]{Definition}

% Patch prob environment to be single spaced
\let\oldprob\prob
\let\endoldprob\endprob
\renewenvironment{prob}
  {\begin{singlespace}\oldprob}
  {\endoldprob\end{singlespace}}

% start problem one line below like for enumerated problems with multiple parts
\newcommand{\belowtitle}{\leavevmode\newline}
%\Observe command
\newcommand{\Observe}{\text{Observe.}}
%(=>)
\newcommand{\IF}{\mathbf{(\Rightarrow)}}
%(<=)
\newcommand{\FI}{\mathbf{(\Leftarrow)}}
%equivalence classes; \class[S]{ *content in square brackets* }
\newcommand{\class}[2][]{\ensuremath{\left[\,#2\,\right]_{#1}}}

\newcommand{\horrule}[1]{\rule{\linewidth}{#1}}
\newcommand{\kkk}{\ensuremath{\Bbbk}}
\newcommand{\CC}{\ensuremath{\mathbb{C}}}
\newcommand{\FF}{\ensuremath{\mathbb{F}}}
\newcommand{\KK}{\ensuremath{\mathbb{K}}}
\newcommand{\NN}{\ensuremath{\mathbb{N}}}
\newcommand{\QQ}{\ensuremath{\mathbb{Q}}}
\newcommand{\RR}{\ensuremath{\mathbb{R}}}
\newcommand{\ZZ}{\ensuremath{\mathbb{Z}}}
\newcommand{\MM}{\ensuremath{\mathcal{M}}}
\newcommand{\TT}{\ensuremath{\mathcal{T}}}
\newcommand{\BB}{\ensuremath{\mathcal{B}}}
\newcommand{\VV}{\ensuremath{\mathcal{V}}}
\newcommand{\WW}{\ensuremath{\mathcal{W}}}
\newcommand{\UU}{\ensuremath{\mathcal{U}}}
\newcommand{\PP}{\ensuremath{\mathcal{P}}}
\newcommand{\LL}{\ensuremath{\mathcal{L}}}
\newcommand{\kk}{\ensuremath{\mathds{k}}}
\newcommand{\EE}{\ensuremath{\mathbb{E}}}
\newcommand{\Ss}{\ensuremath{\mathcal{S}}}
\newcommand{\GG}{\ensuremath{\mathcal{G}}}

\newcommand{\sm}{\char`\\}
%vector stuff
\DeclarePairedDelimiter{\ip}{\langle}{\rangle} %inner product/generate
\DeclarePairedDelimiter{\norm}{\lVert}{\rVert} %norm
\DeclarePairedDelimiter{\sqb}{\lbrack}{\rbrack} %corrd

\newcommand{\floor}[1]{\left\lfloor #1 \right\rfloor}
\newcommand{\ceil}[1]{\left\lceil #1 \right\rceil}
\newcommand{\mbf}[1]{\ensuremath{\mathbf{#1}}}
\newcommand{\tbf}[1]{\textbf{ #1 }}
\newcommand{\Span}{\ensuremath{\mathrm{Span}}}
\newcommand{\Char}[1]{\mathrm{Char}\; #1}
\DeclareMathOperator{\lcm}{lcm}
\newcommand{\id}{\ensuremath{\mathrm{id}}}
\newcommand{\Gal}[2]{\mathrm{Gal}(#1/#2)}
\newcommand{\Aut}{\ensuremath{\mathrm{Aut}}}
\newcommand{\Fix}[2]{\mathrm{Fix}_{#1}(#2)}
\newcommand{\nil}{\ensuremath{\mathrm{nil}}}
\newcommand{\Hom}{\ensuremath{\mathrm{Hom}}}
\newcommand{\Obj}{\ensuremath{\mathrm{Obj}}}

\newcommand{\Rng}{\ensuremath{\mathrm{Rng}}}
\newcommand{\Ring}{\ensuremath{\mathrm{Ring}}}
\newcommand{\h}{\ensuremath{\mathrm{ht}}}


\makeatletter
\renewcommand*\env@matrix[1][*\c@MaxMatrixCols c]{%
  \hskip -\arraycolsep
  \let\@ifnextchar\new@ifnextchar
  \array{#1}}
\makeatother

\def\env@matrix{\hskip -\arraycolsep
  \let\@ifnextchar\new@ifnextchar
  \array{*\c@MaxMatrixCols c}}

\newcommand{\proj}[2]{\text{proj}_{#1}(#2)}

%%% Formatting: Page Header
\newcommand{\StudentName}{Danny Banegas}
\newcommand{\AssignmentName}{Homework 4}
\newcommand{\CourseName}{MATH 720 - Algebra I}

\pagestyle{headandfoot}
\runningheadrule
\firstpageheadrule
\firstpageheader{\CourseName}{\StudentName}{\AssignmentName}
\runningheader{\CourseName}{\StudentName}{\AssignmentName}
\firstpagefooter{}{\thepage}{}
\runningfooter{}{\thepage}{}

\printanswers

\DeclareMathAlphabet{\mathcal}{OMS}{cmsy}{m}{n}

\usepackage{parskip}
\usepackage{setspace}

\begin{document}

%%%%%%%%%%%%%%%%%%%%%%%%%%%%%%%%%%%%%%%%%%%%%%%%% 1 %%%%%%%%%%%%%%%%%%%%%%%%%%%%%%%%%%%%%%%%
\setcounter{thm}{0} % next prob is 1
%%%%%%%%%%%%%%%%%%%%%%%%%%%%%%%%%%%%%%%%%%%%%%%%% lemmas %%%%%%%%%%%%%%%%%%%%%%%%%%%%%%%%%%%%%
\begin{lem}\label{lem:UFDiv}
    If $R$ is a UFD, and $r=\prod_{i=1}^{n}p_{i}$ is some irreducible decomposition, then
    $$d\mid r \iff d=u\prod_{j\in S\subseteq [n]}p_{i}\text{ for some subset }S\subseteq [n]\text{ and some unit }u\in R.$$
\end{lem}

\begin{proof}
    Let $d=d_{1}\cdots d_{m}$ be some irreducible decomposition, which is unique up to associates since $R$ is a UFD. $(\implies)$ If $d\mid r,\;\exists q\in R$ such that $r=dq.$ and there exists some irreducible decomposition $q=q_{1}\cdots q_{t}.$
    
    $r=\prod_{i=1}^{n}p_{i}=dq=(d_{1}\cdots d_{m})(q_{1}\cdots q_{t}).$ Recall that since $R$ is a UFD, and $r$ decomposes into $n$ many $p_{i}$'s, $r=dq\implies dq$ must decompose into $n$ many $d_{i}$'s and $q_{i}$'s, that is $t+m=n.$ So reindex the $q_{i}$'s so that $q_{i}=q_{i+m}$ for all $i\in [t]$ and we obtain $r=p_{1}\cdots p_{m}q_{m+1}\cdots q_{n}.$ Since $R$ is a UFD, there exists a permutation $\sigma\in S_{n}$ such that for some units $u_{1},\hdots, u_{n}\in R$,
    $$d_{i}=u_{i}p_{\sigma(i)},\forall i\in [m]\text{ and }q_{i}=u_{i}p_{\sigma(i)},\forall m<i\leq n.$$
    That is, if $S=\sigma([m])$, since $R$ is commutative, 
    $$d=(\prod_{i=1}^{m}u_{i})(\prod_{j\in S\subseteq [n]}p_{j})=u\prod_{j\in S\subseteq [n]}p_{j}\text{ where }u=\prod_{i=1}^{n}u_{i}.$$
    $(\impliedby)$ On the other hand, if $d=u\prod_{j\in S\subseteq [n]}p_{j}$ for some unit $u\in R$, we have that: 
    $$(u^{-1}\prod_{k\in [n]\setminus S} p_{k})d=(u^{-1}\prod_{k\in [n]\setminus S} p_{k})(u\prod_{j\in S\subseteq [n]}p_{j})=\prod_{i\in [n]}p_{i}=r\implies d\mid r.$$
    
\end{proof}

\begin{lem}
    If $R$ is Noetherian containing a prime ideal $P$, then $R/P$ is also Noetherian. 
\end{lem}
\begin{proof}
    Let $\pi:R\to R/P$ be the canonical homomorphic projection, and $I_{P}$ any ideal of $R/P$. $\pi$ is surjective, so $I=\pi^{-1}(I_{P})\subseteq R$ must be a finitely generated ideal since $R$ is Noetherian. That is, $I=\pi^{-1}(I_{P})=\langle i_{1},\hdots, i_{n}\rangle\implies I_{P}=\pi(I)=\langle \pi(i_{1}),\hdots, \pi(i_{n})\rangle$. Therefore, $R/P$ is Noetherian.

\end{proof}
\begin{lem}
    Let $R$ be a Noetherian UFD, and let $P$ be a prime ideal of $R$. Then there exist some prime elements $q_{1},\hdots, q_{n}\in P$ such that $P=\langle q_{1},\hdots, q_{n}\rangle$.
\end{lem}
\begin{proof}
    Since $R$ is Noetherian, $P$ is finitely generated. So there exist some $a_{1},\hdots, a_{n}\in P$ such that $P=\langle a_{1},\hdots, a_{n}\rangle$. Since $R$ is a UFD, for each $i\in [n]$, $a_{i}=\prod_{j=1}^{m_{i}}p_{(i)_{j}}$ for some irreducible and therefore prime $p_{(i)_{j}}$'s. Well, $\prod_{j=1}^{m_{i}}p_{(i)_{j}}\in P\implies q_{i}=p_{(i)_{k}}\in P$ for some $k\in [m_{i}]$. So then $a_{i}=q_{i}\prod_{\substack{j\in [m_{ij}]\\ p_{(i)_{j}\neq q_{i}}}} p_{(i)_{j}}\implies P=\langle q_{1},\hdots, q_{n}\rangle.$

\end{proof}

\begin{lem}[Problem 7 from HW 3]\label{lem:minprime}
Let $R$ be a commutative ring, and let $I \subseteq R$ be a proper ideal. Prove that the set of prime ideals containing $I$ has minimal elements. We call such a minimal element a \underline{minimal prime ideals over $I$}
\end{lem}

\begin{proof}
    Let $S$ be the set of prime ideals of $R$ that contain $I$. In \textbf{Homework 2} we proved that the set of all proper ideals containing $I$ is a poset with set inclusion. So then $(S\setminus \{R\},\subseteq)$ is an induced subposet of that poset. Obviously, $P\subseteq R,\;\forall P\in S\setminus\{R\}$ and so $(S,\subseteq)$ is also a poset. Now consider any chain $C$ in $S$.

    Let $L_{C}=\cap_{P\in C}P$. Obviously, $L_{C}\subseteq C,\forall P\in C\implies I\subseteq L_{C}$. Next, $\forall a,b\in L_{C}, a,b\in P,\forall P\in C\implies a-b\in P\forall P\in C$ since they are all additive subgroups of $R$. So $L_{C}\leq_{+}R$. Also, $\forall r\in R, \forall a\in L_{C}, ra\in P, \forall P\in C\implies ra=ar\in L_{C}$. So $L_{C}$ is an ideal of $R$ containing $I$. Finally, if $ab\in L_{C},\;$ then $ab\in P,\;\forall P\in C\implies a\text{ or }b\in P,\;\forall P\in C\implies a\text{ or }b\in L_{C}$. So $L_{C}$ is a prime ideal of $R$ containing $I$, and a lower bound for $C$ in $S$. So every chain in $S$ has a lowerbound. 
    
    So then in the dual poset $S^{*}$ every chain has an upperbound in $S$, and by Zorn's Lemma, there exists a maximal element in $S^{*}$. That is, there exists a minimal element in $S$.
    
\end{proof}
\begin{lem}\label{lem:UFDkrull} A prime ideal $P$ of a UFD $R$ has height at most $1$ if and only if $P$ is principal. 
\end{lem}
\begin{proof}
$(\implies)$ If $P$ has height $1$, then there is no prime ideal strictly between $P$ and $\{0\}$. If $P=\{0\}=\langle 0\rangle$, then it's principal. If $P$ is non-zero, suppose it is not principal. So $P$ contains at least two non-zero coprime generators $x$ and $y$. $R$ is a UFD, so $x$ has a prime decomposition $x=p_{1}\cdots p_{n}$. Therefore, since $P$ is prime, some $p_{j}$ belongs to $P$. So $\langle p_{i}\rangle$ is prime and ${0}\subset \langle p_{i}\rangle \subseteq P\implies P$ is principal or $\{0\}\subset \langle p_{j}\rangle \subset P\implies \mathrm{ht}(P)> 1,$ both contradictions. Thus, $P$ is principal.

$(\impliedby)$ If $P=\{0\}=\langle 0\rangle$, it obviously has height $0\leq 1$. If $P$ is principal and prime, then $P=\langle p\rangle$ for some prime element $p\in R\setminus\{0\}$. Suppose there exists some prime ideal $Q$ strictly between $\{0\}$ and $P$. Then there exists some non-zero $q\in Q$ which has a prime decomposition $q=q_{1}\cdots q_{n}$ since $R$ is a UFD and since $Q$ is prime, some $q_{j}$ belongs to $Q$. Recall that all primes are irreducible and vice versa, so then since $q_{j}\in Q\subset P\implies \exists r\in R\text{ s.t. } q_{j}=rp$, $r$ or $p$ is a unit. No primes are units, so $r$ is a unit. Well, then $q_{j}=rp\in Q\implies (p\in Q\implies P=\langle p\rangle \subseteq Q)\text{ or }(r\in Q\implies rr^{-1}=1\in Q\implies Q=R\subset P),$ which are both contradictions. Thus, there is no prime ideal strictly between $\{0\}$ and $P$, so $\mathrm{ht}(P)=1$.
\end{proof}
\newpage
\begin{lem}\label{lem:NoP}
    If $R$ is Noetherian domain, then a principal prime ideal $P$ of $R$ has height at most $1$
\end{lem}
\begin{proof}
Let $R$ be a Noetherian integral domain. If $P=\{0\}=\langle 0\rangle$, it has height $0\leq 1$. If $P$ is prime, principal, and non-zero, then $P=\langle p\rangle$ for some prime element $p\in R\setminus\{0\}$. Suppose there exists a prime ideal $Q$ such that $\{0\}\subset Q\subset P.$ Then $\exists q\in Q\setminus\{0\}\subset P\implies \exists r_{0}\in R$ such that $q=r_{0}p$. Since $Q$ is prime we have that $r_{0}\in Q \text{ or } p\in Q.$  but $p$ can't belong to $Q$ since $Q\subset P=\langle p\rangle$, so $r_{0}\in Q$.

We can define a sequence $(r_{i})_{i\in \NN}$ in $R$ recursively by repeating the same process indefinitely
$$r_{i}p^{i+1}\in Q\implies r_{i}\in Q\implies r_{i}=r_{i+1}p\text{ for some }r_{i+1}\in R\implies q=r_{i}p^{i+1},\;\forall i\in \NN.$$
Now, since $r_{i}=r_{i+1}p,\forall i\in\NN$, we obtain an ascending chain of ideals
$$ 
\langle r_{0}\rangle \subseteq \langle r_{1}\rangle \subseteq \cdots \subseteq \langle r_{i}\rangle \subseteq \cdots
$$
which must stabilize at some $i=n\in \ZZ^{+}$ since $R$ is Noetherian. We already have that $r_{n}=r_{n+1}p$ and now $\langle r_{n}\rangle=\langle r_{n+1}\rangle\implies r_{n+1}=sr_{n}$ for some $s\in R$. So then $r_{n}=r_{n+1}p=(sr_{n})p\implies r_{n}-sr_{n}p=r_{n}(1-sp)=0\implies 1-sp=0.$ But then since $q=sr_{n}p^{n+1}\neq 0$ and $R$ is an integral domain, $s$ is non zero and $p$ is a prime unit. The very thing is impossible.

Therefore there is no prime ideal strictly between $\{0\}$ and $P$, and $\operatorname{ht}(P)\leq 1.$

\end{proof}
\begin{cor}[Corollary of Homework 3 Problem 9]\label{cor:localprime}
    If $R$ is a commutative ring and $P$ is a prime ideal of $R$, then the set of all prime ideals of $R$ contained in $P$ is in bijection with the set of all prime ideals of $R_{P}$ via $Q\longleftrightarrow QR_{P}$.
\end{cor}
\begin{proof}
    Since $P$ is prime and $R$ is commutative, $R_{P}=\{\frac{r}{q}\mid r\in R, q\in R\setminus P\}$ is a local ring with maximal ideal $PR_{P}$. Additionally, by \textbf{Homework 3}, all ideals of $R_{P}$ are of the form $IR_{P}=\{\frac{i}{q}\in R_{P}\mid i\in I\}$ for some $I\subseteq R$, and $I\not\subseteq P\implies IR_{P}=R_{P}$. We show an ideal $Q\subset P$ is prime if and only if $QR_{P}$ is prime.
    
    $(\implies)$ Suppose $Q\subset P$ is prime. If $(\frac{a}{r})(\frac{b}{s})\in QR_{P}$, then $\frac{ab}{rs}\in QR_{P}\implies ab\in Q\implies a\in Q\text{ or }b\in Q\implies \frac{a}{q}\in QR_{P}\text{ or }\frac{b}{d}\in QR_{P}$. So $Q\subset R$ is prime $\implies QR_{P}$ is prime. 
    
    $(\impliedby)$ On the other hand, if $QR_{P}$ is prime, $ab\in Q\implies \frac{ab}{1}=(\frac{a}{1})(\frac{b}{1})\in QR_{P}\implies \frac{a}{1}\in QR_{P}\text{ or }\frac{b}{1}\in QR_{P}\implies a\in Q\text{ or }b\in Q.$ So $QR_{P}$ is prime $\implies Q\subset P$ is prime.

    So we see that $QR_{P}$ is prime $\iff Q\subset P$ is prime and immediately the mapping $\phi$ from prime ideals $Q$ properly contained in $P$ to prime ideals of $R_{P}$ defined by $\phi(Q)=QR_{P}$ is a surjection. Now, if $Q_{1}\neq Q_{2}$ but both are prime ideals properly contained in $P$, consider $Q_{1}R_{P}$ and $Q_{2}R_{P}$. Since $Q_{1}\neq Q_{2}$, there exists some $q\in Q_{1}\setminus Q_{2}$ or some $q\in Q_{2}\setminus Q_{1}$. Without loss of generality, let $q\in Q_{1}\setminus Q_{2}$. Then $\frac{q}{1}\in Q_{1}R_{P}$ but $\frac{q}{1}\not\in Q_{2}R_{P}$, so $Q_{1}R_{P}\neq Q_{2}R_{P}$. Therefore, $\phi$ is injective. So $\phi$ is a bijection.

    \textbf{(2)} Now consider $(\PP,\subseteq)$ and $(\PP_{P},\subseteq)$ where $\PP$ is the set of all prime ideals properly contained in $P$, and $\PP_{P}$ is the set of all prime ideals in $R_{P}$. All primes in $R$ must contain $\mathrm{nil}(R)$ since $R$ is commutative, and so by \textbf{Homework 4} $(\PP,\subseteq)\subseteq ()$ the induced subposet obtained by removing $\hat{0}=\nil(R)$, every proper ideal is contained in some maximal element,  $\PP$ is  We proved that both of these are posets with $\hat{0}$. Well, if $Q_{1}\subseteq $
\end{proof}
\newpage
\begin{lem}[Krull's Height Theorem]
Let $R$ be Noetherian. If an ideal $I$ is generated by $n$ elements and $P$ is a prime ideal minimal over $I$, then
\[
\operatorname{ht}(P)\leq n.
\]
\end{lem}
\begin{proof}
(a) Let
\[
P_0\supseteq P_1\supseteq P_2\supseteq\cdots
\]
be a descending chain of prime ideals.

Since $R$ is Noetherian, $P_0$ is finitely generated, say
\[
P_0=\langle a_1,\hdots,a_n\rangle.
\]
Since $P_0$ is prime and is minimal over itself, Krull's Height Theorem gives
\[
\operatorname{ht}(P_0)\leq n.
\]
Thus $P_0$ has finite height.

If the descending chain did not stabilize, then after removing repeated terms we would get a strictly descending chain
\[
P_0\supsetneq P_{i_1}\supsetneq P_{i_2}\supsetneq\cdots.
\]
Then for every $m$ we get a strict chain of prime ideals
\[
P_{i_m}\subsetneq P_{i_{m-1}}\subsetneq\cdots\subsetneq P_{i_1}\subsetneq P_0
\]
ending at $P_0$ of length $m$. Since $m$ is arbitrary, this contradicts the fact that $P_0$ has finite height.

Therefore the chain stabilizes.

(b) Let $P$ be a prime ideal of height $1$. Since $P\neq (0)$, choose $0\neq x\in P$. Since $R$ is a UFD,
\[
x=up_1\cdots p_n
\]
where $u$ is a unit and each $p_i$ is irreducible. Since $P$ is prime and $x\in P$, some $p_i\in P$.

In a UFD, irreducibles are prime, so $(p_i)$ is a nonzero prime ideal. Hence
\[
(0)\subsetneq (p_i)\subseteq P.
\]
Since $\operatorname{ht}(P)=1$, we must have
\[
(p_i)=P.
\]
Thus $P$ is principal.
\end{proof}

%%%%%%%%%%%%%%%%%%%%%%%%%%%%%%%%%%%%%%%%%%%%%%%%% 1 %%%%%%%%%%%%%%%%%%%%%%%%%%%%%%%%%%%%%%%%
\newpage
\setcounter{thm}{0} % next prob is 1
\begin{prob}
Let $R$ be a UFD, and let $a$ and $b$ be nonzero non-unit elements of $R$.
Prove that $a$ and $b$ have a gcd, and describe it in terms of their prime factorizations.
\end{prob}

\begin{proof}
    Since $R$ is a UFD, there exist irreducible decompositions $a=a_{1}\cdots a_{n},b=b_{1}\cdots b_{m}.$ Recall that any irreducibles in a UFD are also prime elements, so these are also prime products. Now let $G=\{p_{1},\hdots,p_{N}\}$ be some minimal collection of non-associate primes appearing in both $a$ and $b$, which must exist since both have finite decompositions. Minimal here means we can express both $a$ and $b$ as a product within $G$. Then there exists $\alpha_{i}$ and $\beta_{i}$ such that $a=\prod_{i=1}^{N} g_{i}^{\alpha_{i}}$ and $b=\prod_{i=1}^{N} g_{i}^{\beta_{i}}$. 
    
    Now let $d=\prod_{i=1}^{N} g_{i}^{\gamma_{i}}$ where $\gamma_{i}=\min\{\alpha_{i},\beta_{i}\}$ for each $i\in [N]$. Behold, $\gamma_{i}\leq \alpha_{i},\beta_{i}$ and so
    \begin{align*}
        g\prod_{i=1}^{N} g_{i}^{\alpha_{i}-\gamma_{i}}=(\prod_{i=1}^{N} g_{i}^{\alpha_{i}})(\prod_{i=1}^{N} g_{i}^{\alpha_{i}-\gamma_{i}})=a\implies g\mid a\\
        g\prod_{i=1}^{N} g_{i}^{\beta_{i}-\gamma_{i}}=(\prod_{i=1}^{N} g_{i}^{\beta_{i}})(\prod_{i=1}^{N} g_{i}^{\beta_{i}-\gamma_{i}})=b\implies g\mid b
    \end{align*}
 Now consider any common divisor $d$ of $a$ and $b$. By \textbf{Lemma }\ref{lem:UFDiv}, $d$ can be expressed as some product of primes appearing in $a$ and $b$ and some units. Then obviously it is expressible as some unit multiplied by some product in $G$. That is, $d=u\prod_{i=1}^{N} g_{i}^{\delta_{i}}$ for some unit $u\in R$ and $\delta_{i}\leq \alpha_{i},\beta_{i}$ for each $i\in [N]$ since these are non associate primes. Well, $\gamma_{i}=\min\{\alpha_{i},\beta_{i}\}\in \{\alpha_{i},\beta_{i}\}$ for each $i\in [N]$, so $\delta_{i}\leq \alpha_{i},\beta_{i}\implies \delta_{i}\leq \gamma_{i},\;\forall i\in [N]$. Therefore,
    $$d\prod_{i=1}^{N} g_{i}^{\gamma_{i}-\delta_{i}}=(\prod_{i=1}^{N} g_{i}^{\delta_{i}})(\prod_{i=1}^{N} g_{i}^{\gamma_{i}-\delta_{i}})=\prod_{i=1}^{N} g_{i}^{\gamma_{i}}=g\implies d\mid g.$$
Thus, $g$ is a gcd of $a$ and $b$.
   
\end{proof}

%%%%%%%%%%%%%%%%%%%%%%%%%%%%%%%%%%%%%%%%%%%%%%%%% 2 %%%%%%%%%%%%%%%%%%%%%%%%%%%%%%%%%%%%%%%%
\newpage
\begin{prob}
Let $R$ be a PID and let $S$ be a multiplicative set of $R$, and let $P \subseteq R$ be a prime ideal.
\begin{enumerate}[label=(\alph*)]
\item Prove that $R/P$ is a PID.
\item Prove that $S^{-1}R$ is a PID.
\end{enumerate}
\end{prob}
\begin{proof}
    \textbf{(a)} Consider any ideal $I_{P}\subseteq R/P$ and let $\pi:R\to R/P$ be the canonical homomorphic projection. Then the premiage $I=\pi^{-1}(I_{P})$ of $I_{P}$ is an ideal. Since $R$ is a PID, $I=\langle i\rangle$ for some $i \in R$. Then $I_{P}=\pi(I)=\pi(\langle i\rangle)=\langle \pi(i)\rangle$, so $I_{P}$ is principal.

    \textbf{(b)} Consider any ideal $I_{S}\subseteq S^{-1}R$ and let $\iota:R\to S^{-1}R$ be the inclusion mapping defined by $\iota(r)=\frac{r}{1}$, which is clearly a homomorphism. Then the premiage $I=\iota^{-1}(I_{S})$ of $I_{S}$ is an ideal. Since $R$ is a PID, $I=\langle i\rangle$ for some $i \in R$. Then $I_{S}=\iota(I)=\iota(\langle i\rangle)=\langle \iota(i)\rangle$, so $I_{S}$ is principal.
\end{proof}

%%%%%%%%%%%%%%%%%%%%%%%%%%%%%%%%%%%%%%%%%%%%%%%%% 3 %%%%%%%%%%%%%%%%%%%%%%%%%%%%%%%%%%%%%%%%
\newpage
\begin{prob}
Let $R$ be a Noetherian UFD.
\begin{enumerate}[label=(\alph*)]
\item Prove that every descending chain of prime ideals
$$
P_0 \supseteq P_1 \supseteq \cdots \supseteq P_n \supseteq \cdots
$$
in $R$ stabilizes.
\item The height of a prime ideal $P$ is (if finite) the maximum length $h$ of a chain of prime ideals
$$
P_0 \subseteq P_1 \subseteq \cdots \subseteq P_h = P.
$$
Prove that every prime ideal in $R$ of height $1$ is principal.
\end{enumerate}
\end{prob}


\begin{proof}\textbf{(b)} This is a direct consequence of \textbf{Lemma }\ref{lem:UFDkrull}. If $P$ has height $1$, then $P$ is principal.

\textbf{(a)} By some lemmas, every prime ideal $P_{i}=\langle p_{i,1},\hdots, p_{i,n_{i}}\rangle$ is such that each generator $p_{i,j}$ has a prime decomposition unique up to associates so $p_{i,j}=\prod q_{i,j,k}$ and then by induction some $q_{i,j,k*}$ must belong to $P_{i}$ since $P_{i}$ is prime. Relabel each such $q_{i,j,k*}$ as $q_{i,j}$ and then we have that $P_{i}=\langle q_{i,1},\hdots, q_{i,n_{i}}\rangle$ is a finite prime generating presentation for each $P_{i}$ in our chain.

Next we need to use a minimal prime ideal over $P_{i}'s$ but I need a lot of localization machinery to essentially prove Krull's Principal Ideal theorem...


\end{proof}

%%%%%%%%%%%%%%%%%%%%%%%%%%%%%%%%%%%%%%%%%%%%%%%%% 4 %%%%%%%%%%%%%%%%%%%%%%%%%%%%%%%%%%%%%%%%
\newpage
\begin{prob}
Consider the ring
$$
R := \ZZ[\sqrt{-5}] := \{a + b\sqrt{-5} \mid a,b \in \ZZ\}.
$$
\begin{enumerate}[label=(\alph*)]
\item Prove that $R \cong \ZZ[x]/(x^2+5)$.
\item Since $R$ is a subring of $\QQ(\sqrt{d})$ where $d=-5$, the norm $N$ defined in Homework 2 will still have the property that
$$
N(zw)=N(z)N(w)
$$
in $R$. Use this to prove that the units in $R$ are $\pm 1$.
\item Prove that $2$, $3$, $(1+\sqrt{-5})$, and $(1-\sqrt{-5})$ are all irreducible non-associate elements of $R$.
\item Prove that $R$ is not a UFD.
\end{enumerate}
\end{prob}

%%%%%%%%%%%%%%%%%%%%%%%%%%%%%%%%%%%%%%%%%%%%%%%%% 5 %%%%%%%%%%%%%%%%%%%%%%%%%%%%%%%%%%%%%%%%
\newpage
\begin{prob}
Let $R$ be a Euclidean domain which is not a field. Prove that there exists a nonzero, non-unit element $c \in R$ such that for all $a \in R$, there is some $q,r \in R$ such that
$$
a = qc + r
$$
and $r=0$ or $r$ is a unit.

\tbf{Hint:} Pick $c$ so that $N(c)$ is a minimal positive value among norms of non-units.
\end{prob}

%%%%%%%%%%%%%%%%%%%%%%%%%%%%%%%%%%%%%%%%%%%%%%%%% 6 %%%%%%%%%%%%%%%%%%%%%%%%%%%%%%%%%%%%%%%%
\newpage
\begin{prob}
Consider the ring
$$
R := \ZZ\left[\frac{1+\sqrt{-19}}{2}\right]
:= \left\{a + b\frac{1+\sqrt{-19}}{2} \,\middle|\, a,b \in \ZZ\right\}.
$$
\begin{enumerate}[label=(\alph*)]
\item Since $R$ is a subring of $\QQ(\sqrt{d})$ where $d=-19$, the norm $N$ defined in Homework 2 will still have the property that
$$
N(zw)=N(z)N(w)
$$
in $R$. Prove that the smallest values of $N(z)$ for $z\in R$ are $0,1,4,5$ and that
$$
N\!\left(a + b\frac{1+\sqrt{-19}}{2}\right) \ge 5 \quad \text{if } b\ne 0.
$$
\item Prove that the units in $R$ are $\pm 1$.
\item Prove that if $c \in R$ satisfies the condition in the previous exercise, then $c$ must divide $2$ or $3$ in $R$, and conclude $c=\pm 2$ or $\pm 3$.
\item Show that there is no $q,r$ such that
$$
\frac{1+\sqrt{-19}}{2}=qc+r
$$
with $c=\pm 2,\pm 3$ and $r\in\{0,\pm 1\}$, and conclude that $R$ is not a Euclidean domain.
\end{enumerate}
\end{prob}

\end{document}